\documentclass[11pt,reqno]{amsart}

\usepackage[T1]{fontenc}
\usepackage{lmodern}
\usepackage{amsmath,amssymb,mathtools}
\usepackage{xcolor}
\usepackage{microtype}
\usepackage[colorlinks=true,linkcolor=blue!55!black,citecolor=blue!55!black,urlcolor=blue!55!black]{hyperref}
\hypersetup{
  pdftitle={Localized Stationary Profiles in a Singular Diffusion Limit of the Brusselator},
  pdfauthor={Haibo Lu},
  pdfkeywords={Brusselator, localized stationary pattern, reversible spatial dynamics, singular diffusion, homoclinic orbit}
}

\numberwithin{equation}{section}

\newtheorem{theorem}{Theorem}[section]
\newtheorem{lemma}[theorem]{Lemma}
\newtheorem{proposition}[theorem]{Proposition}
\newtheorem{corollary}[theorem]{Corollary}
\theoremstyle{remark}
\newtheorem{remark}[theorem]{Remark}

\newcommand{\R}{\mathbb R}
\newcommand{\Fix}{\operatorname{Fix}}
\newcommand{\cR}{\mathcal R}
\newcommand{\cM}{\mathcal M}
\newcommand{\cO}{\mathcal O}
\newcommand{\norm}[1]{\left\lVert #1\right\rVert}

\title[Localized Brusselator profiles]{Localized Stationary Profiles in a Singular Diffusion Limit of the Brusselator}
\author{Haibo Lu}
\address{Haibo Lu}
\email{luhaibo1985@gmail.com}
\date{August 28, 2026}

\subjclass[2020]{35B36, 35K57, 34C37, 34E15, 37C29}
\keywords{Brusselator, localized stationary pattern, reversible spatial dynamics, singular diffusion, homoclinic orbit}

\begin{document}

\begin{abstract}
We consider the classical one-dimensional Brusselator at fixed feed parameters
$A=B=1$, with activator diffusivity $d>0$ and inhibitor diffusivity one.
Assuming the computer-assisted Core Lemma stated below, we prove that for every
sufficiently small $d>0$ the homogeneous state $(1,1)$ possesses a nonconstant
even localized stationary profile with positive concentrations.  The proof
rescales the joint limit $d\to0^+$ to a fixed reversible saddle-focus core,
continues one transverse intersection of its unstable manifold with the
reverser fixed set, and reflects the resulting half-orbit.  The continued
profiles are uniformly exponentially localized and satisfy
\[
 \begin{gathered}
 \norm{u_d-1}_{\infty}=\Theta(d^{1/2}),\qquad
 \norm{v_d-1}_{\infty}=\Theta(d),\\
 W_u(d),W_v(d)=\Theta(d^{1/4}),
 \end{gathered}
\]
where $W_u,W_v$ are central connected half-height widths.  The result is a
stationary existence theorem; it does not assert temporal stability, dynamical
selection from the Turing instability, or a canard interpretation of the
localized orbit.
\end{abstract}

\maketitle

\section{Introduction}

The Brusselator is a standard activator--inhibitor model for diffusion-driven
pattern formation.  We study the one-dimensional system
\begin{equation}\label{eq:pde}
 \begin{aligned}
 u_t&=d u_{xx}-2u+1+u^2v,\\
 v_t&=v_{xx}+u-u^2v,
 \end{aligned}
 \qquad x\in\R,
\end{equation}
whose homogeneous state is $(u,v)=(1,1)$.  Our parameter is the activator
diffusivity $d$, and the inhibitor diffusivity has been normalized to one.
The question is whether a particular homoclinic geometry at the singular
limit persists for all sufficiently small positive $d$ and yields a genuine
positive-concentration stationary profile of \eqref{eq:pde}.

Localized Brusselator states are not new.  Reversible $1{:}1$ resonance gives
small homoclinic solutions near an ordinary Turing point
\cite{IoossPeroueme1993,GlebskyLerman1997}; spike, mesa, snaking, and
two-dimensional localized regimes have also been studied analytically and
numerically
\cite{TzouEtAl2011,TzouNecWard2013,KolokolnikovErneuxWei2006,AlSaadiEtAl2021,VillarChampneys2023,KostetEtAl2018}.
On a bounded domain, computer-assisted branches of stationary and periodic
Brusselator solutions and Hopf bifurcations have been proved in
\cite{Arioli2022}.
The singular pulses in \cite{TzouEtAl2011} instead use a finite interval,
boundary inhibitor feed, global activator feed, and simultaneously small feed
parameters; their strong activator spikes and $\cO(1)$ inhibitor scale do not
give the fixed-$A=B=1$ whole-line branch studied here.
The regime here has an additional singular quantifier.  If
$\varepsilon=\sqrt d$, the Turing value for the general Brusselator at
$A=1$ is
\begin{equation}\label{eq:turing-path}
 B_T=(1+\varepsilon)^2,
 \qquad B-B_T=-2\varepsilon-\varepsilon^2
 \quad\text{when }B=1.
\end{equation}
Thus $d\to0^+$ follows a prescribed path that approaches the Turing curve.
The physical critical wave number $k_T=d^{-1/4}$ diverges; equivalently, the
wave number $\varepsilon k_T=d^{1/4}$ in the fast variable $y=x/\varepsilon$
tends to zero.  The normal-form coefficients also become singular.

For each fixed sufficiently small $\varepsilon>0$, Proposition~2.1 of
\cite{JencksEtAl2026} supplies local reversible-Turing homoclinics for $B<B_T$
and $B$ close enough to $B_T$, in the appropriate subcriticality interval.
That statement has the quantifier pattern
\[
 \text{$\varepsilon$ fixed}\quad\Longrightarrow\quad
 \text{there is a neighborhood $\delta(\varepsilon)>0$ of $B_T$}.
\]
It does not state a lower bound
$\delta(\varepsilon)>2\varepsilon+\varepsilon^2$, which would be needed to
substitute the path \eqref{eq:turing-path}.  The contribution of this paper
is to close precisely this joint-limit gap for one selected orbit and to
translate the continuation into uniform physical scale laws.

Set
\begin{equation}\label{eq:r-def}
 r=d^{1/4}=\sqrt\varepsilon.
\end{equation}
An exact weighted scaling transforms the stationary equation into a smooth
polynomial family whose $r=0$ member is a fixed reversible Hamiltonian
saddle-focus system.  A computer-assisted calculation in a separate frozen
source provides one symmetric homoclinic of that core and, crucially, a
nonzero two-dimensional shooting determinant.  We state this input as the
Core Lemma in Section~\ref{sec:core}.  Everything after that lemma---positive
parameter invariant manifolds, reversible matching, full-line tails,
positivity, and physical scale estimates---is proved analytically here.

The proof has five steps.  Section~\ref{sec:scaling} derives the exact
positive-$r$ spatial system.  Section~\ref{sec:core} states the sole
model-specific imported input and explains its transversality content.
Section~\ref{sec:manifolds} proves parameter-uniform local stable and unstable
manifolds on fixed weighted spaces.  Section~\ref{sec:matching} continues the
core intersection and constructs a global symmetric homoclinic without using
a positive-$r$ first integral.  Section~\ref{sec:pde-conclusions} returns to
the PDE and proves positivity, amplitudes, exponential localization, and
widths.  Section~\ref{sec:scope} separates these conclusions from stability,
Turing selection, and canard questions.

\section{Exact central scaling and the main result}\label{sec:scaling}

The stationary equation associated with the general Brusselator is
\begin{equation}\label{eq:general-stationary}
 \begin{aligned}
 0&=d u_{xx}-(B+1)u+A+u^2v,\\
 0&=v_{xx}+Bu-u^2v.
 \end{aligned}
\end{equation}
Put $\varepsilon=\sqrt d$, $p=\varepsilon u_x$, $q=v_x$, and
$y=x/\varepsilon$.  Then \eqref{eq:general-stationary} is equivalent to
\begin{equation}\label{eq:fast-system}
 \begin{aligned}
 u_y&=p,& p_y&=-A+(B+1)u-u^2v,\\
 v_y&=\varepsilon q,& q_y&=\varepsilon(-Bu+u^2v).
 \end{aligned}
\end{equation}
Both the physical-$x$ and fast-$y$ systems are reversed by
\begin{equation}\label{eq:physical-reverser}
 \cR(u,p,v,q)=(u,-p,v,-q).
\end{equation}

We now fix $A=B=1$, use \eqref{eq:r-def}, introduce
$\xi=ry=x/r$, and set
\begin{equation}\label{eq:central-scaling}
 u=1+r^2U,\qquad p=r^3P,\qquad
 v=1+r^4V,\qquad q=r^3Q.
\end{equation}
A direct expansion, with no remainder, gives the polynomial family
\begin{equation}\label{eq:Fr}
 \begin{aligned}
 U'&=P,\\
 P'&=-U^2-V-2r^2UV-r^4U^2V,\\
 V'&=Q,\\
 Q'&=U+r^2(U^2+V)+2r^4UV+r^6U^2V,
 \end{aligned}
\end{equation}
where the prime denotes $d/d\xi$.  We write its vector field as $F_r$.
It is reversible under
\begin{equation}\label{eq:central-reverser}
 \cR(U,P,V,Q)=(U,-P,V,-Q),
 \qquad \Fix\cR=\{P=Q=0\}.
\end{equation}
At $r=0$, \eqref{eq:Fr} reduces exactly to
\begin{equation}\label{eq:core-system}
 U'=P,\qquad P'=-U^2-V,\qquad V'=Q,\qquad Q'=U.
\end{equation}
The limiting system has primitive and Hamiltonian
\begin{equation}\label{eq:core-hamiltonian}
 \lambda_0=P\,dU-Q\,dV,
 \qquad
 H_0=\frac12(Q^2-P^2)-\frac13U^3-UV,
\end{equation}
with $\iota_{F_0}d\lambda_0=dH_0$.  We use
\eqref{eq:core-hamiltonian} only to identify the imported core orbit; no
Hamiltonian structure is asserted for $r>0$.

The linearization $A_r=DF_r(0)$ has characteristic polynomial
\begin{equation}\label{eq:charpoly}
 \mu^4-r^2\mu^2+1
\end{equation}
and, for $0\le r<\sqrt2$, the saddle-focus quartet
\begin{equation}\label{eq:eigenvalues}
 \pm\alpha(r)\pm i\beta(r),
 \qquad
 \alpha(r)=\frac12\sqrt{2+r^2},\quad
 \beta(r)=\frac12\sqrt{2-r^2}.
\end{equation}
In particular, stable and unstable dimensions are both two and the spectral
gap is uniform near $r=0$.

For $0<\eta<2^{-1/2}$, define
\begin{equation}\label{eq:Xeta}
 X_\eta=\left\{Z\in C^0(\R,\R^4):
 \norm{Z}_\eta:=\sup_{\xi\in\R}e^{\eta|\xi|}|Z(\xi)|<\infty\right\}.
\end{equation}
For an even stationary profile, define its central connected half-height
widths by
\begin{equation}\label{eq:width-def}
 \begin{aligned}
 W_u(d)&=\operatorname{length}\,\mathcal C_0
 \left\{x:|u_d(x)-1|\ge\frac12|u_d(0)-1|\right\},\\
 W_v(d)&=\operatorname{length}\,\mathcal C_0
 \left\{x:|v_d(x)-1|\ge\frac12|v_d(0)-1|\right\},
 \end{aligned}
\end{equation}
where $\mathcal C_0$ denotes the connected component containing zero.

\begin{theorem}[Localized positive Brusselator profiles]\label{thm:main}
Assume the computer-assisted Core Lemma~\ref{lem:core}.  There exist
$r_0>0$, $C>0$, and $0<\eta<2^{-1/2}$, together with a $C^1$ map
\[
 [0,r_0]\longrightarrow X_\eta,\qquad r\longmapsto Z_r,
\]
where differentiability is with respect to $r=d^{1/4}$, such that:
\begin{enumerate}
\item $Z_0$ is the selected core orbit from Lemma~\ref{lem:core}.  For
every $0<r\le r_0$, $Z_r$ is a nonconstant $\cR$-symmetric homoclinic
orbit of \eqref{eq:Fr}, centered at the locally continued intersection
$c_r=Z_r(0)\in\Fix\cR$, and $c_r\to c_0$ as $r\to0^+$.
\item The tails are uniform:
\begin{equation}\label{eq:uniform-tail}
 |Z_r(\xi)|\le Ce^{-\eta|\xi|}
 \quad(0\le r\le r_0,\ \xi\in\R),
 \qquad \norm{Z_r-Z_0}_\eta\longrightarrow0
 \quad\text{as }r\to0^+.
\end{equation}
\item If $d=r^4$, then
\begin{equation}\label{eq:inverse-scaling}
 u_d(x)=1+r^2U_r(x/r),\qquad
 v_d(x)=1+r^4V_r(x/r)
\end{equation}
is a smooth, nonconstant, even stationary solution of \eqref{eq:pde}, with
$u_d(x)>0$, $v_d(x)>0$, and $(u_d(x),v_d(x))\to(1,1)$ as
$|x|\to\infty$.
\item The amplitudes satisfy
\begin{equation}\label{eq:amplitude-limits}
 \begin{aligned}
 d^{-1/2}\norm{u_d-1}_\infty&\longrightarrow\norm{U_0}_\infty>0,\\
 d^{-1}\norm{v_d-1}_\infty&\longrightarrow\norm{V_0}_\infty>0.
 \end{aligned}
 \qquad(d\to0^+)
\end{equation}
\item There are $0<c<C_w<\infty$, independent of $d$, such that
\begin{equation}\label{eq:physical-localization}
 d^{-1/2}|u_d(x)-1|+d^{-1}|v_d(x)-1|
 \le Ce^{-\eta|x|/d^{1/4}}
 \qquad(x\in\R,\ 0<d\le r_0^4)
\end{equation}
and
\begin{equation}\label{eq:width-law}
 c d^{1/4}\le W_u(d),W_v(d)\le C_wd^{1/4}
 \qquad(0<d\le r_0^4).
\end{equation}
\end{enumerate}
The threshold $d_0=r_0^4$ is existential; it is not an explicitly certified
diffusion interval.
\end{theorem}

\section{The computer-assisted core input}\label{sec:core}

We isolate the sole model-specific imported statement.  This formulation is
deliberately smaller than the return--exit and coding results available in the
source: those additional conclusions are not used here.

\begin{lemma}[Computer-assisted Core Lemma]\label{lem:core}
The core system \eqref{eq:core-system} has a selected nonconstant
$\cR$-symmetric homoclinic orbit $\Gamma_0$ to the origin.  Its center
$c_0=\Gamma_0(0)$ lies in $\Fix\cR$ and satisfies
\begin{equation}\label{eq:center-enclosures}
 \begin{aligned}
 U_0(0)&\in[4.8785234574459304,4.8785234988116768],\\
 V_0(0)&\in[-7.9333304994224827,-7.933330385013492].
 \end{aligned}
\end{equation}
Moreover,
\begin{equation}\label{eq:core-transversality}
 W^u_0(0)\pitchfork\Fix\cR
 \quad\text{at }c_0.
\end{equation}
More precisely, for the source phase $\phi$ and finite flight time $T$, the
normal shooting map
\[
 \cM_0(\phi,T)=(P,Q)\bigl(\Phi_0^T(S_0(\phi))\bigr)
\]
has a locally unique zero $(\phi_0,T_0)$ in the reported shooting box and
\begin{equation}\label{eq:det-enclosure}
 \det D_{(\phi,T)}\cM_0(\phi_0,T_0)
 \in[149.56393055300413,149.56404227745782].
\end{equation}
The corresponding hit is the first positive nontrivial hit of $\Fix\cR$.
\end{lemma}

The lemma is imported from Proposition~8.6 together with the frozen
certificate and validation README
described in Appendix~\ref{app:evidence}.  The manuscript's equations
(8.36)--(8.40) give the phase, flight time, determinant, tangent rows, and
first-hit identification.  The center enclosures, local uniqueness in the
shooting box, and detailed sign checks excluding an earlier hit are recorded
in the certificate and README.  The determinant has the following direct
geometric meaning.  Let
\[
 G(\phi,T)=\Phi_0^T(S_0(\phi)),
 \qquad N(U,P,V,Q)=(P,Q).
\]
Then $\cM_0=N\circ G$, $\ker DN=T\Fix\cR$, and the image of $DG$ is the
tangent space of the flowed two-dimensional unstable patch.  Invertibility
of $D(N\circ G)$ implies that this tangent space intersects $T\Fix\cR$
trivially, which is exactly \eqref{eq:core-transversality}.  This is
transversality to the reverser fixed set, not ambient transversality of
$W^u$ and $W^s$ in four dimensions.

\begin{remark}\label{rem:evidence-status}
Lemma~\ref{lem:core} is a computer-assisted imported input, not a floating
point observation.  Its frozen source records outward-rounded flow and first
variation bounds, a two-dimensional Krawczyk inclusion, a rigorous local
unstable graph, and first-hit sign checks.  At the date of this manuscript,
the immutable source revision and hashes are recorded, but the complete
evidence closure is not anonymously accessible at a stable public archive.
Thus Theorem~\ref{thm:main} is mathematically conditional on the precisely
stated Core Lemma, and a publication release must also make that lemma's
evidence package externally auditable.
\end{remark}

\section{Uniform local invariant manifolds}\label{sec:manifolds}

We prove the parameter-uniform local result needed both for continuation and
for full-line tails.  Write
\begin{equation}\label{eq:linear-nonlinear}
 F_r(Z)=A_rZ+N_r(Z),
 \qquad N_r(0)=DN_r(0)=0.
\end{equation}
Although the physical problem has $r\ge0$, the polynomial family extends to
negative $r$, which permits an ordinary $C^1$ parameter argument at zero.

\begin{lemma}[Uniform local manifolds]\label{lem:uniform-manifolds}
After decreasing $r_0>0$, the origin of $F_r$ has two-dimensional local
stable and unstable manifolds $W^s_{\mathrm{loc}}(r)$ and
$W^u_{\mathrm{loc}}(r)$ for $|r|\le r_0$.  They admit graph and half-orbit
parameterizations over fixed two-dimensional spaces which are $C^1$ in the
graph parameter and in $r$.  The half-orbits and their first $r$-derivatives
obey exponential estimates with constants independent of $r$.  A compact
patch may be flowed for every bounded time for which its $r=0$ flowout lies
in a compact subset of the flow domain, and the continued flowout retains
this $C^1$ dependence.
\end{lemma}

\begin{proof}
Choose $r_0<1$.  By \eqref{eq:eigenvalues}, the stable and unstable Riesz
projections $\Pi_r^s,\Pi_r^u$ vary smoothly.  Fix
\begin{equation}\label{eq:rates}
 0<\eta<\eta_1<a_0<\inf_{|r|\le r_0}\alpha(r).
\end{equation}
After decreasing $r_0$, the restrictions
\begin{equation}\label{eq:fixed-spaces}
 J_r^u=\Pi_r^u|_{E_0^u}:E_0^u\to E_r^u,
 \qquad
 J_r^s=\Pi_r^s|_{E_0^s}:E_0^s\to E_r^s
\end{equation}
are smooth isomorphisms, and their inverses are uniformly bounded.  Spectral
calculus in finite dimension and the Duhamel formula give $K_j$ such that,
for $j=0,1$,
\begin{equation}\label{eq:semigroup-estimates}
 \begin{aligned}
 \norm{\partial_r^j(e^{A_rt}\Pi_r^s)}
 &\le K_j(1+t)^je^{-a_0t} &&(t\ge0),\\
 \norm{\partial_r^j(e^{A_rt}\Pi_r^u)}
 &\le K_j(1+|t|)^je^{a_0t} &&(t\le0).
 \end{aligned}
\end{equation}

For the unstable manifold use the fixed space
\begin{equation}\label{eq:minus-space}
 X^-_{\eta_1}=\left\{Z\in C^0(({-}\infty,0],\R^4):
 \norm{Z}_{\eta_1,-}=\sup_{t\le0}e^{-\eta_1t}|Z(t)|<\infty\right\}.
\end{equation}
Apply an $\cR$-invariant radial cutoff at radius $2\delta$ to $N_r$, equal
to $N_r$ on the ball of radius $\delta$, and call the result
$\widetilde N_r$.  Its uniform Lipschitz constant $L_\delta$ tends to zero
with $\delta$.  For $b\in E_0^u$ define
\begin{equation}\label{eq:LP-map}
 \begin{aligned}
 (\mathcal T_{r,b}Z)(t)={}&e^{A_rt}J_r^ub
 +\int_0^t e^{A_r(t-s)}\Pi_r^u\widetilde N_r(Z(s))\,ds\\
 &+\int_{-\infty}^t e^{A_r(t-s)}\Pi_r^s
       \widetilde N_r(Z(s))\,ds.
 \end{aligned}
\end{equation}
With
\[
 D_0=(a_0-\eta_1)^{-1}+(a_0+\eta_1)^{-1},
 \qquad K=\max\{K_0,1\},
\]
the estimates above imply
\begin{equation}\label{eq:contraction-estimates}
 \begin{aligned}
 \norm{\mathcal T_{r,b}Z}_{\eta_1,-}
 &\le K\norm{J_r^u}|b|+KL_\delta D_0\delta,\\
 \norm{\mathcal T_{r,b}Z-\mathcal T_{r,b}\widetilde Z}_{\eta_1,-}
 &\le KL_\delta D_0\norm{Z-\widetilde Z}_{\eta_1,-}.
 \end{aligned}
\end{equation}
Choose $\delta$ so that $KL_\delta D_0\le1/2$, and then choose $b_*>0$
so that
\[
 K\sup_{|r|\le r_0}\norm{J_r^u}b_*\le\delta/2.
\]
Thus \eqref{eq:LP-map} is a uniform strict contraction of one fixed closed
ball for all $|r|\le r_0$ and $|b|\le b_*$.  Denote its fixed point by
$z^u(r,b)$.

The strict gaps in \eqref{eq:rates} also control the differentiated kernels:
the factor $(1+|t-s|)e^{-a_0|t-s|}$ remains integrable in the
$\eta_1$-weighted norm, and $\partial_r\widetilde N_r(Z)=\cO(|Z|^2)$
uniformly in the cutoff ball.  The parameterized contraction theorem gives
\begin{equation}\label{eq:C1-half-orbit}
 (r,b)\longmapsto z^u(r,b)\in X^-_{\eta_1}
 \quad\text{of class }C^1.
\end{equation}
The endpoint
\begin{equation}\label{eq:Gu}
 G^u(r,b)=z^u(r,b)(0)
\end{equation}
parameterizes $W^u_{\mathrm{loc}}(r)$.  The fixed point lies in the region
where the cutoff agrees with the original field.  The stable construction is
the forward-time analogue, and reversibility plus local uniqueness gives
\begin{equation}\label{eq:stable-reversibility}
 W^s_{\mathrm{loc}}(r)=\cR W^u_{\mathrm{loc}}(r).
\end{equation}
The final assertion follows from the common-flow-domain theorem and smooth
finite-time dependence for ODEs on compact subsets of the flow domain.
\end{proof}

\section{Reversible matching and the full homoclinic}\label{sec:matching}

Let $S_0(\phi)$ be the source curve used in Lemma~\ref{lem:core}, with
matching map $\cM_0$.  Choose a compact phase interval $I$ containing
$\phi_0$ in its interior.  Since $S_0(I)$ is a compact curve in
$W^u_0(0)$, after shrinking $I$ there is a common $\tau>0$ for which
\begin{equation}\label{eq:pull-source-back}
 \widehat S_0(\phi)=\Phi_0^{-\tau}(S_0(\phi))
 =G^u(0,b_0(\phi)),
 \qquad |b_0(\phi)|<b_*,
\end{equation}
for a $C^1$ map $b_0:I\to E_0^u$.  Define the genuine
parameter-dependent unstable source curve by
\begin{equation}\label{eq:continued-source}
 \widehat S(r,\phi)=G^u(r,b_0(\phi)),
 \qquad
 S(r,\phi)=\Phi_r^\tau(\widehat S(r,\phi)).
\end{equation}
Then $S(r,\phi)\in W^u_r(0)$, $S(0,\phi)=S_0(\phi)$, and the family is
$C^1$.

Choose a compact time interval $J\Subset(0,\infty)$ containing $T_0$.
After shrinking $I,J$, all base trajectories
$\Phi_0^T(S_0(\phi))$, $(\phi,T)\in I\times J$, lie in a common compact
flow tube.  Therefore
\begin{equation}\label{eq:matching-map}
 \cM(r,\phi,T)=(P,Q)\bigl(\Phi_r^T(S(r,\phi))\bigr)
\end{equation}
is $C^1$ near $(0,\phi_0,T_0)$.  Lemma~\ref{lem:core} and the
finite-dimensional implicit-function theorem give $C^1$ functions
$\phi(r),T(r)$ with
\begin{equation}\label{eq:implicit-branch}
 \cM(r,\phi(r),T(r))=0,
 \qquad \phi(0)=\phi_0,\quad T(0)=T_0.
\end{equation}
After decreasing $r_0$, continuity of the determinant, together with the
base enclosure, gives, for example,
\begin{equation}\label{eq:continued-det-bound}
 \left|\det D_{(\phi,T)}\cM(r,\phi(r),T(r))\right|\ge70
 \qquad(|r|\le r_0).
\end{equation}
This numerical constant is not a certified positive-$r$ interval statement;
it is a convenient consequence of continuity on a sufficiently small
existential neighborhood.

Put
\begin{equation}\label{eq:center-cr}
 c_r=\Phi_r^{T(r)}(S(r,\phi(r))).
\end{equation}
Then $c_r\in\Fix\cR$, $c_r\in W^u_r(0)$, and $c_r\ne0$ for small $|r|$.
The negative half-orbit exists for all negative time and converges to zero.
Define the positive half by reflection:
\begin{equation}\label{eq:reflected-orbit}
 Z_r(\xi)=
 \begin{cases}
  \Phi_r^\xi(c_r),&\xi\le0,\\
  \cR\Phi_r^{-\xi}(c_r),&\xi\ge0.
 \end{cases}
\end{equation}
Because $c_r\in\Fix\cR$ and $D\cR F_r=-F_r\circ\cR$, the values and
first derivatives agree at zero.  ODE uniqueness makes
\eqref{eq:reflected-orbit} a smooth, nonconstant, symmetric global
homoclinic.  Notice that forward global existence was constructed by
reflection rather than inferred from a nonexistent positive-$r$ energy
integral.

We next record the full-line weighted dependence.  The entry-to-center time
$\sigma(r)=\tau+T(r)$ is bounded.  Choose a fixed
$T_*>\sup_{|r|\le r_0}\sigma(r)$ large enough that
\begin{equation}\label{eq:ar-definition}
 \begin{aligned}
 a_r&:=z^u\!\left(r,b_0(\phi(r))\right)
       \left(-(T_*-\sigma(r))\right)\\
    &=Z_r(-T_*)
 \end{aligned}
\end{equation}
lies strictly in the uniform graph disk.  Formula
\eqref{eq:ar-definition} and \eqref{eq:C1-half-orbit} show directly that
$r\mapsto a_r$ is $C^1$; no full-line Banach-space regularity has been used.
Set
\begin{equation}\label{eq:br-definition}
 b_r=(J_r^u)^{-1}\Pi_r^u a_r.
\end{equation}
Then $b_r$ is $C^1$ and $|b_r|<b_*$.  Evaluating
\eqref{eq:LP-map} at $t=0$ gives
\begin{equation}\label{eq:LP-coordinate}
 \Pi_r^uG^u(r,b)=J_r^ub.
\end{equation}
Since $a_r$ lies on the local unstable graph,
\eqref{eq:br-definition}--\eqref{eq:LP-coordinate} and uniqueness of the
graph parameter imply $a_r=G^u(r,b_r)$.  Orbit uniqueness now gives the
fixed-shift representation
\begin{equation}\label{eq:fixed-shift-tail}
 Z_r(\xi)=z^u(r,b_r)(\xi+T_*)
 \qquad(\xi\le-T_*).
\end{equation}
Equations \eqref{eq:C1-half-orbit}, \eqref{eq:fixed-shift-tail}, and
$\eta<\eta_1$ yield
\begin{equation}\label{eq:negative-tail-C1}
 \sup_{|r|\le r_0}\sup_{\xi\le-T_*}e^{-\eta\xi}
 \bigl(|Z_r(\xi)|+|\partial_rZ_r(\xi)|\bigr)<\infty.
\end{equation}
Common finite-time dependence controls $[-T_*,0]$, and reflection controls
the positive half-line.  Hence
\begin{equation}\label{eq:full-C1}
 r\longmapsto Z_r\in X_\eta\quad\text{is }C^1,
 \qquad
 \sup_{|r|\le r_0}\sup_{\xi\in\R}e^{\eta|\xi|}
 \bigl(|Z_r(\xi)|+|\partial_rZ_r(\xi)|\bigr)<\infty.
\end{equation}
This proves the first two assertions of Theorem~\ref{thm:main}.

\section{Return to the PDE: positivity and scale laws}\label{sec:pde-conclusions}

The inverse change of variables \eqref{eq:inverse-scaling} is exact for every
$r>0$, so $Z_r$ gives a stationary solution of \eqref{eq:pde}.  Symmetry of
$Z_r$ gives
\[
 U_r(-\xi)=U_r(\xi),\qquad V_r(-\xi)=V_r(\xi),
\]
and hence the physical profiles are even.

From \eqref{eq:full-C1}, there is $M>0$ such that
\begin{equation}\label{eq:uniform-components}
 \sup_{0\le r\le r_0}
 \bigl(\norm{U_r}_\infty+\norm{V_r}_\infty\bigr)\le M.
\end{equation}
Decrease $r_0$ so that $r_0^2M<1/2$ and $r_0^4M<1/2$.  Then
\begin{equation}\label{eq:positive}
 u_d(x)\ge\frac12,
 \qquad v_d(x)\ge\frac12
\end{equation}
for every $x\in\R$ and $0<r\le r_0$.  Exponential convergence to the
homogeneous state follows directly from \eqref{eq:uniform-tail}.

Uniform convergence $Z_r\to Z_0$ and the exact identities
\begin{equation}\label{eq:exact-amplitudes}
 d^{-1/2}\norm{u_d-1}_\infty=\norm{U_r}_\infty,
 \qquad
 d^{-1}\norm{v_d-1}_\infty=\norm{V_r}_\infty
\end{equation}
prove \eqref{eq:amplitude-limits}.  Both limits are nonzero by
\eqref{eq:center-enclosures}.  Replacing $\xi$ by $x/r$ in
\eqref{eq:uniform-tail}, and using $r=d^{1/4}$, proves
\eqref{eq:physical-localization}.

It remains to prove the width statement for the explicit observable
\eqref{eq:width-def}.  By \eqref{eq:center-enclosures} and uniform
convergence there is $m>0$ with
\begin{equation}\label{eq:center-lower}
 |U_r(0)|\ge m,
 \qquad |V_r(0)|\ge m
 \qquad(0\le r\le r_0).
\end{equation}
The weighted bound on $P_r=U_r'$ and $Q_r=V_r'$ gives a uniform derivative
bound $M_1\ge1$.  With $\ell=m/(4M_1)$, the mean-value theorem gives
\begin{equation}\label{eq:inner-width}
 |U_r(\xi)|\ge\frac12|U_r(0)|,
 \qquad |V_r(\xi)|\ge\frac12|V_r(0)|
 \qquad(|\xi|\le\ell).
\end{equation}
Choose $L>0$, independent of $r$, so that $Ce^{-\eta L}<m/2$.
The two scaled half-height inequalities fail for $|\xi|\ge L$, while
\eqref{eq:inner-width} shows that their central components contain
$[-\ell,\ell]$.  Their lengths in $\xi$ lie between $2\ell$ and $2L$.
Multiplication by $x=r\xi$ proves \eqref{eq:width-law}, and completes the
proof of Theorem~\ref{thm:main}.

\section{Scope and relation to Turing patterns and canards}\label{sec:scope}

Theorem~\ref{thm:main} proves the existence of a spatially localized
stationary solution.  It also determines how its amplitude and one precise
width observable scale along the singular path.  These are physical-variable
consequences of the homoclinic continuation, but they do not by themselves
describe what the time-dependent PDE selects from generic initial data.

Three nearby questions require additional arguments.

First, temporal stability is a spectral problem for the linearization of
\eqref{eq:pde} about $(u_d,v_d)$.  The instability theorem in
\cite{GlebskyLerman1997} concerns specified branches in a fixed-diffusion
reversible-Hopf setting; its hypotheses and singular uniformity would have to
be checked before either stability or instability could be transferred to
the branch constructed here.

Second, \eqref{eq:turing-path} locates the parameter path relative to the
Turing curve, but the present proof starts from a selected RFSN-II core and
continues its spatial homoclinic.  It does not prove that a branch selected by
time evolution or numerical continuation from the Turing bifurcation is the
same global branch for every small $d$.

Third, spatially periodic canards organize important Brusselator dynamics in
the same singular regime \cite{JencksEtAl2026}.  The local oscillations of the
core used here arise from its saddle-focus eigenvalues, whereas a canard
designation requires a verified orbit segment that tracks attracting and
repelling slow manifolds through the relevant charts.  No such identification
is needed for Theorem~\ref{thm:main}, and none is claimed.

These distinctions leave a useful hierarchy.  The present theorem supplies
a rigorously organized stationary branch, scale laws, and a canonical profile
on which later numerical and spectral studies can be based.  Explicit
positive-$d$ validation, temporal stability, Turing-branch identification,
and canard geometry are subsequent problems rather than omitted consequences
of the theorem proved here.

\appendix

\section{Frozen evidence provenance}\label{app:evidence}

The Core Lemma is imported from H.~Lu, \emph{First returns, singular exits,
and action finite parts near a reversible Hamiltonian saddle-focus}, at the
frozen revision of
\url{https://github.com/h-lu/reversible-rfsn-ii-waves},
\begin{center}
\texttt{d54add098545063d5efe8f1d6f062d4cfc116a0d}.
\end{center}
The cited manuscript location is Proposition~8.6, specifically its
verification of the transverse homoclinic tube.  Equations (8.36)--(8.40)
record the phase, flight time, determinant, tangent rows, and first-hit
identification, including
\begin{equation}\label{eq:phase-time-enclosures}
 \begin{aligned}
 \phi_0&\in[5.8615055856447817,5.8615055856450482],\\
 T_0&\in[9.6374420678958099,9.6374420678971511].
 \end{aligned}
\end{equation}
The endpoint enclosures \eqref{eq:center-enclosures}, local uniqueness in the
shooting box, and detailed early-hit sign exclusions are recorded in the
frozen certificate and its README rather than in those displayed manuscript
equations.
The frozen SHA-256 identities are as follows:
\begin{center}
\small
source manuscript\\[-2pt]
\texttt{0baf6335aad72d5893479d8876d2613671ecb8ac2ccd73664405dea4381e6a20}\\[3pt]
validation README\\[-2pt]
\texttt{34f3e15f475d22be939020bde2b8480ae33ba60a9c043f577088ab5d70253610}\\[3pt]
certificate JSON\\[-2pt]
\texttt{ed0f9f58f8ba5f1d5c36dc7c3a72bb725599c4172a3cd610d890b88699fecfbd}\\[3pt]
interval source\\[-2pt]
\texttt{655cfe16e2cdb24185f11cbb314e7c0b2f029705d4588fdcbaf5a9f0993298f3}.
\end{center}
The interval calculation used outward-rounded CAPD/FILIB flow and
first-variation enclosures, a two-dimensional Krawczyk inclusion, and
interval sign checks for the first symmetry hit.  Its rigorous local unstable
graph is a required dependency of the shooting certificate, so the
certificate JSON alone is not a complete evidence archive.

The source URL and historical hashes preserve provenance.  A submission
release must additionally provide anonymous stable access either to the
frozen source and its complete evidence closure or to an equivalent archived
bundle with replay metadata.  This packaging obligation affects external
auditability, not the analytic implication from Lemma~\ref{lem:core} to
Theorem~\ref{thm:main}.

\bibliographystyle{abbrv}
\bibliography{references}

\end{document}
